%!TEX root =  ../main.tex
\renewcommand{\columnseprule}{1.5pt}
\begin{multicols*}{2}
\rule[0.5\baselineskip]{0.4\textwidth}{1pt}
\noindent
\LabSection{Truth to Power}\label{sec:0504p}
\begin{exercises}{sec:0504p}
\lab{} Making every box equal to 0.5, graph of $f(x)=x^{1/3}$ below.  Also graph $g(x)=x^1$.  One what intervals is the first greater than the second?

\noindent
\begin{tikzpicture}[xscale=0.8,yscale=0.8]
	\draw [thick, <->] (-5.3,0) -- (5.3,0);
	\draw [thick, ->] (0,-5.3) -- (0,5.3);
	\draw [thin] (-5,-5) grid (5,5);
\end{tikzpicture}

\vspace{3cm}
\lab{} On the same graph, also graph $h(x)=x^{5/3}$.  On what intervals is this greater than $g(x)$?

\vspace{3cm}
\lab{} When are the functions increasing and decreasing?

\subprob $f(x)$
\subprob $g(x)$
\subprob $h(x)$

\lab{}  Record the derivative of all three functions.



\vspace{4cm}
\lab{}  When are the derivatives positive and negative?
\subprob $f^\prime(x)$
\subprob $g^\prime(x)$
\subprob $h^\prime(x)$

\lab{}  Find the derivatives of the derivatives, the second derivatives of each function.  Record them here.


\vspace{4cm}
\lab{} When are the function concave and convex?  Concave is ‘U’ shaped, while convex is hill-shaped.
\subprob $f(x)$
\subprob $g(x)$
\subprob $h(x)$

\lab{} When are the second derivatives positive and negative?
\subprob $f^{\prime\prime}(x)$
\subprob $g^{\prime\prime}(x)$
\subprob $h^{\prime\prime}(x)$

\lab{} In technical language, describe what you think the point of this problem set is.


\end{exercises}
\end{multicols*}